\documentclass[9pt,landscape,a4paper]{article}
\usepackage[utf8]{inputenc}
\usepackage[T1]{fontenc}
\usepackage{geometry}
\usepackage{multicol}
\usepackage{amsmath,amssymb}
\usepackage{enumitem}
\usepackage{graphicx}
\usepackage{xcolor}
\usepackage{listings}
\usepackage{booktabs}
\usepackage{titlesec}
\usepackage{fancyhdr}
\usepackage{hyperref}
\usepackage{CJKutf8}

% Page layout
\geometry{top=0.7cm, bottom=0.7cm, left=0.7cm, right=0.7cm}
\pagestyle{empty}
\setlength{\parindent}{0pt}
\setlength{\parskip}{1.5pt}
\setlength{\columnsep}{7pt}

% Section formatting
\titleformat{\section}{\normalsize\bfseries\color{blue!70!black}}{}{0em}{}[\titlerule]
\titleformat{\subsection}{\small\bfseries\color{blue!50!black}}{}{0em}{}
\titlespacing{\section}{0pt}{4pt}{2pt}
\titlespacing{\subsection}{0pt}{3pt}{1pt}

% Code listing style
\lstset{
    language=Python,
    basicstyle=\ttfamily\fontsize{6.5}{7.5}\selectfont,
    keywordstyle=\color{blue!70!black}\bfseries,
    stringstyle=\color{red!60!black},
    commentstyle=\color{green!50!black}\itshape,
    numberstyle=\tiny\color{gray},
    backgroundcolor=\color{gray!8},
    frame=single,
    framerule=0.3pt,
    rulecolor=\color{gray!50},
    breaklines=true,
    breakatwhitespace=false,
    tabsize=2,
    showstringspaces=false,
    aboveskip=2pt,
    belowskip=2pt,
    xleftmargin=2pt,
    xrightmargin=2pt,
}

\newcommand{\keybox}[1]{%
    \begin{center}
    \fcolorbox{blue!50!black}{blue!5}{%
        \parbox{0.95\linewidth}{\small #1}%
    }
    \end{center}
}

\newcommand{\zhint}[1]{{\scriptsize\color{gray!70!black}#1}}

\begin{document}
\begin{CJK}{UTF8}{gbsn}

\begin{center}
{\large\bfseries\color{blue!70!black} Chapter 10 -- Case Studies (II): Natural Language Problems}\\
{\footnotesize IERG4320/IEMS5726 Data Science in Practice | Cheat Sheet}
\end{center}

\begin{multicols}{3}

% ==================== SECTION 1 ====================
\section{NLP Overview}
\zhint{NLP：让计算机理解和生成自然语言}

\textbf{Applications}: sentiment analysis, text classification, chatbots, machine translation, summarization, speech recognition, image captioning.\\
\textbf{Why NLP is hard}: high-dimensional, sparse, sequential \& context-dependent, noisy \& unstructured, language-specific.\\
\zhint{高维、稀疏、上下文依赖、噪声多、语言差异}

\textbf{LLM}: trained on vast text via self-supervised learning; billions--trillions of parameters. Transformer architecture (2017) enabled GPT, BERT, LLaMA, DeepSeek, Gemini.

\subsection{Semantic Analysis}
\zhint{语义分析：从文本中提取情感/上下文/含义}
\begin{itemize}[leftmargin=8pt,itemsep=0pt,topsep=0pt]
\item \textbf{Context}: word meaning study (translation)
\item \textbf{Emotions}: happy, sad, fear, calm
\item \textbf{Sentiments}: positive, neutral, negative
\end{itemize}

\textbf{Non-ML concepts}: hyponyms, meronomy, polysemy, synonyms, antonyms, homonyms.\\
\zhint{上下位词、部分整体、一词多义、近义词、反义词、同形异义}

\subsection{Classic ML Approach}
\zhint{经典方法：TF-IDF + 朴素贝叶斯，准确率约0.84}\\
TF-IDF features + Naive Bayes classifier on IMDB (25k train / 25k test). Accuracy $\approx$ 0.84.\\
\textbf{Limitations}: cannot handle polysemy/homonyms, no word position, simple model.

% ==================== SECTION 2 ====================
\section{Neural Bag-of-Words (NBoW)}
\zhint{NBoW：词嵌入取平均后送入全连接网络分类}

\textbf{Steps}: tokenize $\to$ embedding lookup $\to$ average pooling $\to$ FC layers $\to$ classification.\\
\zhint{分词→查嵌入→平均池化→全连接层→分类}

\begin{lstlisting}
class NeuralBagOfWords(nn.Module):
  def __init__(self, vocab_size, embed_dim,
               hidden_dim, pad_idx):
    super().__init__()
    self.embedding = nn.Embedding(
      vocab_size, embed_dim, padding_idx=pad_idx)
    self.fc1 = nn.Linear(embed_dim, hidden_dim)
    self.fc2 = nn.Linear(hidden_dim, hidden_dim)
    self.fc3 = nn.Linear(hidden_dim, 1)
    self.dropout = nn.Dropout(0.3)
    self.relu = nn.ReLU()
    self.sigmoid = nn.Sigmoid()
  def forward(self, x):
    embedded = self.embedding(x)
    bow = torch.mean(embedded, dim=1) # avg pool
    out = self.dropout(self.relu(self.fc1(bow)))
    out = self.dropout(self.relu(self.fc2(out)))
    return self.sigmoid(self.fc3(out)).squeeze(-1)
# Result: Test Acc = 0.8639 (5 epochs)
\end{lstlisting}

\textbf{Pros}: simple, fast, parallelizable, fixed-size output.\\
\textbf{Cons}: ignores word order, no context, limited semantics.\\
\zhint{优：简单快速; 缺：忽略词序和上下文}

% ==================== SECTION 3 ====================
\section{LSTM for NLP}
\zhint{LSTM：按时间步顺序处理，保留上下文记忆}

Sequential processing: each word is a time step.\\
Takes current embedding + previous hidden state $\to$ gates decide what to remember/forget.\\
\zhint{每个词按顺序处理，门控机制决定记忆/遗忘}

\textbf{Pros}: captures word order, long-range dependencies.\\
\textbf{Cons}: sequential (slow), fails on very long sequences.\\
Test Acc $\approx$ 0.53 with only 5 epochs (needs longer training).

% ==================== SECTION 4 ====================
\section{BERT (Bidirectional Encoder)}
\zhint{BERT：双向Transformer编码器，Google 2018}

\textbf{Key ideas}: bidirectional context, self-attention, pre-train then fine-tune.\\
\zhint{双向上下文 + 自注意力 + 预训练后微调}

\subsection{Input Representation}
\texttt{input = token\_emb + segment\_emb + position\_emb}\\
Special tokens: \texttt{[CLS]} (classification), \texttt{[SEP]} (separator).\\
Tokenizer: WordPiece (subword-level).

\subsection{Pre-training Strategies}
\begin{enumerate}[leftmargin=10pt,itemsep=0pt,topsep=0pt]
\item \textbf{MLM} (Masked Language Model): mask 15\% tokens, predict originals \zhint{遮蔽语言模型}
\item \textbf{NSP} (Next Sentence Prediction): predict if sentence B follows A \zhint{下一句预测}
\end{enumerate}

\subsection{BERT Variants}
\zhint{BERT变体：更小更快的版本}
\begin{itemize}[leftmargin=8pt,itemsep=0pt,topsep=0pt]
\item \textbf{ALBERT}: lighter, less training time
\item \textbf{RoBERTa}: Facebook, improved pre-training
\item \textbf{ELECTRA}: replaced token detection
\item \textbf{DistilBERT}: compact, faster (66M params)
\item \textbf{TinyBERT}: similar to DistilBERT
\end{itemize}

\subsection{Code: DistilBERT Mini}
\begin{lstlisting}
config = {'vocab_size': 30522, 'dim': 256,
  'n_heads': 4, 'n_layers': 4,
  'intermediate_dim': 1024, 'dropout': 0.1,
  'max_length': 256, 'num_labels': 2}

class DistilBertMini(nn.Module):
  def __init__(self, config):
    super().__init__()
    self.embeddings = nn.Embedding(
      config['vocab_size'], config['dim'])
    encoder_layer = nn.TransformerEncoderLayer(
      d_model=config['dim'],
      nhead=config['n_heads'],
      dim_feedforward=config['intermediate_dim'],
      dropout=config['dropout'],
      batch_first=True, norm_first=True)
    self.transformer = nn.TransformerEncoder(
      encoder_layer,
      num_layers=config['n_layers'])
    self.layer_norm = nn.LayerNorm(config['dim'])
    self.dropout = nn.Dropout(config['dropout'])
    self.classifier = nn.Linear(
      config['dim'], config['num_labels'])
  def forward(self, input_ids, attention_mask):
    x = self.embeddings(input_ids)
    key_pad = ~attention_mask.bool()
    x = self.transformer(
      x, src_key_padding_mask=key_pad)
    # Mean pooling over non-padded tokens
    masked = x * attention_mask.unsqueeze(-1)
    pooled = masked.sum(1) / attention_mask\
      .sum(1, keepdim=True)
    pooled = self.dropout(self.layer_norm(pooled))
    return self.classifier(pooled)
# Params: 10.9M (16.6% of DistilBERT)
# Test Acc: 0.8371 (5 epochs)
\end{lstlisting}

\textbf{Applications}: semantic similarity, text classification, QA, summarization, NER.

% ==================== SECTION 5 ====================
\section{Transformer Model}
\zhint{Transformer：完全基于自注意力，无需RNN/CNN}

\textbf{Self-attention}: captures interdependencies between all words simultaneously (not sequential).\\
\textbf{3 attention types}: within encoder, encoder--decoder, within decoder.\\
\zhint{三种注意力：编码器内、编解码器间、解码器内}

\textbf{Always use pre-trained models} -- training from scratch is very expensive (BERT: 8 days on TPU v3; re-training: 2--3 hours).
\begin{lstlisting}
# TinyBERT pre-trained transformer
from transformers import (AutoTokenizer,
  AutoModelForSequenceClassification)
name = "huawei-noah/TinyBERT_General_4L_312D"
tokenizer = AutoTokenizer.from_pretrained(name)
model = AutoModelForSequenceClassification\
  .from_pretrained(name, num_labels=2)
# Params: 14.4M, Test Acc: 0.8566 (3 epochs)
\end{lstlisting}

\subsection{Encoder vs Decoder vs Seq2Seq}
\zhint{编码器(BERT) vs 解码器(GPT) vs 编解码器(T5)}
\begin{itemize}[leftmargin=8pt,itemsep=0pt,topsep=0pt]
\item \textbf{Encoder-only} (BERT): bidirectional, strong feature extraction, not for generation
\item \textbf{Decoder-only} (GPT): autoregressive, great for text generation, unidirectional
\item \textbf{Seq2Seq} (T5): encoder--decoder, complex I/O, most expensive
\end{itemize}

% ==================== SECTION 6 ====================
\section{GLUE \& SQuAD Benchmarks}
\zhint{GLUE：9个NLP任务的标准评估基准}

\textbf{GLUE} (9 tasks):
\begin{itemize}[leftmargin=8pt,itemsep=0pt,topsep=0pt]
\item \textbf{CoLA}: grammatical correctness (Matthews Corr.)
\item \textbf{SST-2}: sentiment classification (Accuracy)
\item \textbf{MRPC}: paraphrase detection (F1/Acc)
\item \textbf{QQP}: question duplicate detection (F1/Acc)
\item \textbf{STS-B}: semantic similarity (Pearson/Spearman)
\item \textbf{MNLI}: natural language inference (Accuracy)
\item \textbf{QNLI/RTE/WNLI}: inference tasks (Accuracy)
\end{itemize}

\textbf{SQuAD}: reading comprehension.\\
v1.1: extract answer span from paragraph.\\
v2.0: includes unanswerable questions.\\
Metrics: Exact Match (EM) + F1 Score.

% ==================== SECTION 7 ====================
\section{Text Generation}
\zhint{文本生成：AI生成连贯自然的文字，本质是"下一词预测"}

Harder than NLP: must be coherent, different each time, output $>$ input.

\subsection{Markov Chain (Non-DL)}
\zhint{马尔可夫链：根据前一个词的概率预测下一个词}\\
$P(X_n|X_1,...,X_{n-1}) = P(X_n|X_{n-1})$\\
Transition matrix P: row $i$, col $j$ = $P(X_{n+1}=j|X_n=i)$.\\
Steady state: $\pi_n = \pi_{n-1} P$ (converges if irreducible \& aperiodic).

\textbf{Markovify}: Python package for text generation.\\
\textbf{Pros}: no DL needed, fast, light on memory.\\
\textbf{Cons}: mixes styles, no backward prediction, cannot handle polysemy/homonyms.

\subsection{GPT (Decoder-only Transformer)}
\zhint{GPT：用解码器生成类人文本，也叫大语言模型(LLM)}

\begin{itemize}[leftmargin=8pt,itemsep=0pt,topsep=0pt]
\item \textbf{GPT-1} (2018): 117M params, BooksCorpus
\item \textbf{GPT-2} (2019): 1.5B params, WebText (45TB)
\item \textbf{GPT-3} (2020): 175B params, few/zero-shot
\item \textbf{GPT-4} (2023): 1.76T params, multimodal
\end{itemize}

\textbf{Why bigger}: early layers learn word order/POS; middle layers learn semantics; late layers learn high-level abstractions.

\begin{lstlisting}
from transformers import (AutoTokenizer,
  AutoModelForCausalLM)
model_name = "distilgpt2"
tokenizer = AutoTokenizer.from_pretrained(
  model_name)
tokenizer.pad_token = tokenizer.eos_token
model = AutoModelForCausalLM.from_pretrained(
  model_name, dtype=torch.float32,
  device_map="cpu")
\end{lstlisting}

\subsection{LLaMA (Meta)}
\zhint{LLaMA：Meta开源的大语言模型系列}\\
Free for research \& commercial use.\\
v1 (Feb 2023) $\to$ v2 (Jul 2023) $\to$ v3 (Apr 2024) $\to$ v4 (Apr 2025).\\
LLaMA 4 Maverick: 17B active $\times$ 128 experts $\approx$ 400B total.\\
\textbf{Llama Stack}: unified API for inference, RAG, agents, tools, safety, evals.

\keybox{
\textbf{Model Comparison on IMDB}:\\
Naive Bayes + TF-IDF: 0.84 | NBoW: 0.86\\
DistilBERT Mini: 0.84 | TinyBERT: 0.86\\
\textbf{Takeaway}: Transformers win with pre-training + fine-tuning; NBoW is a strong simple baseline.\\
\zhint{Transformer靠预训练+微调取胜; NBoW是简单但有效的基线}
}

\end{multicols}
\end{CJK}
\end{document}
